\section{R8: meaningful results and equivalence}
\label{sec:ranking}

\Q{R8.1}{What ranking objective is defensible?}
\ans A declared, task-conditioned preference model is defensible; a universal numerical notion of ``the intended program'' is not. Use correctness as an admission condition and expose the main preference dimensions, even if the interface keeps the default policy automatic.

A Bayesian or minimum-description-length interpretation is useful. For a noiseless specification $D$, every consistent candidate has the same likelihood under the hard consistency model. Ranking then comes from a prior:
\[
  -\log P(p\mid T,D,E)
  =-\log P(p\mid T,E)+\text{constant},\qquad p\models D.
\]
Thus ``shortest'' is a particular prior, not a consequence of the examples. A useful description length may include the program, invented helper definitions, carrier type, and the cost of invoking rare components. Arbitrarily making one library constant free can reverse the ranking, so the grammar and its version belong in the receipt.

Use hard profile constraints first, then a modest Pareto frontier over implementation size, proof/axiom requirements, and a declared operational or readability proxy. If a single default is needed, define a stable lexicographic or scalar key and state it precisely. The smallest proof need not accompany the simplest program, and the smallest program may be asymptotically slower. These choices are preferences, not logical facts.

Do not make input usage an absolute criterion. Constant functions, projections, and lawful interfaces with redundant parameters legitimately ignore inputs. A syntactic occurrence can also be dead code rather than useful dependence. Treat unused-input count as one low-confidence feature, conditioned on the signature and contract, rather than as evidence of an incorrect or meaningless candidate. A contract saying $f(0)=0$ cannot identify the intended behavior at other inputs; the tool should preserve that ambiguity instead of presenting its prior as a proof of intent.

The inspected engine currently sorts by unused outer binders, eliminator occurrences, and expression size, and deduplicates candidates by pretty-printed text \cite{engine}. These are clear places to separate semantic identity, display preference, and discovery order. A stable final tie-breaker should use a canonical serialized expression and dependency identity, not local timestamps or a pretty-printer's options.

An anytime stop criterion needs the same cost everywhere. If $p^*$ is an incumbent and $L(s)$ is a proved lower bound on the final ranking cost of every completion of frontier state $s$, then
\[
  \min_{s\in\mathrm{frontier}}L(s)\geq c(p^*)
\]
certifies that no strictly cheaper candidate remains, provided the frontier covers all remaining alternatives. It does not certify a richer top-$k$ diversity objective, and it does not apply when the search queue is ranked by an unrelated learned probability. For $k$ distinct semantic classes, the frontier bound and duplicate policy must be defined over those classes as well.

\Q{R8.2}{What inexpensive equivalence is safe for dependent programs?}
\ans Use a hierarchy, and make the level explicit. Canonical syntax, definitional equality, proved extensional equality, and agreement on finitely many observations are different relations with different permissible uses.

Canonical alpha-renaming and safe local reductions provide cheap exact keys. For two closed terms at the same target type, a successful definitional-equality check provides stronger evidence, but a resource-limited failure means unknown. A proved equality may justify merging terms extensionally even when conversion does not. For function types, that may require a function-extensionality proof; its axioms must respect the chosen profile.

A fingerprint of outputs on a finite test set is only an observational cluster. To see why, fix tested natural inputs $0,\ldots,b$ and compare
\[
 f(n)=n,\qquad
 g(n)=\begin{cases}n&n\leq b,\\n+1&n>b.\end{cases}
\]
They have identical fingerprints and differ at $b+1$. The same phenomenon occurs for lists, including shape-parametric functions that change their positional action after a length threshold. Finite observations do not prove equivalence.

For \emph{closed accepted results}, clustering is useful display compression: show one representative, retain alternatives, and record the tests that define the cluster. New probes can split it later. For \emph{partial search states}, deleting every but one fingerprint representative is sound only if the relation is a congruence for every remaining grammar context and the associated proof/interface constraints. That requires a fragment-specific theorem; it is not implied by matching current examples.

Dependent results add a typing requirement. Values at $F(i)$ and $F(j)$ cannot be compared by an ordinary equality until their indices have been related and a transport chosen. Include index and family identity in fingerprints, or normalize them with certified equality evidence. For a dependent function $(x:X)\to F(x)$, compare outputs pointwise in each common input fiber. Comparing arbitrary serialized outputs after erasing their indices is not a semantic equivalence test.

Finally, distinguish extensional equality from implementation equivalence under a cost model. The quadratic reverse fold and a linear accumulator reverse can be extensionally equal while materially different as implementations. A user asking for useful synthesized code may prefer retaining both until the ranking policy accounts for that difference.

\section{R9: model-guided proposals with independent certification}
\label{sec:learning}

\Q{R9.1}{What accuracy do current models achieve on carrier and helper proposals?}
\ans The inspected sources do not establish an accuracy figure for the exact Leant2 task: proposing a useful typed carrier or helper interface, without supplying the whole implementation, under the proposal's component restrictions. General proof, code-generation, or premise-selection scores cannot answer that question. A numerical answer would require an experiment; none was run here.

Define the event before measuring it. A carrier proposal is \emph{well formed} if it elaborates in the allowed grammar. It is \emph{expressive} for a task if some allowed update and finishing program can meet the contract. It is \emph{useful at budget $B$} if the fixed downstream synthesizer actually finds and certifies such a program within that budget. Those probabilities can differ dramatically. Counting a plausible-looking product type as a successful answer measures the first, easiest event at best.

Run four paired conditions with the same downstream engine: blind carrier enumeration; distinction-guided symbolic proposals; a fixed recorded model proposal stream; and an oracle carrier supplied by a verified reference solution. Use identical component environments and either equal total time or clearly separated synthesis-only and end-to-end time. Include model initialization, inference, malformed outputs, and rejected proposals in end-to-end cost.

Report top-1 and top-$k$ well-formedness, certified end-to-end success, useful-proposal latency, and the residual gap to the oracle condition. Stratify by first-order state, function carriers, indexed motives, and helper lemmas. Use disjoint algorithm families for held-out testing, keep prompts and model versions fixed, and publish every proposal rather than only successful ones. Confidence intervals should be computed on independent task units or with family-level resampling when tasks share structure; many nearly identical list variants are not many independent observations.

A successful model-assisted run proves that particular proposal useful under the recorded configuration. It does not prove that learning is necessary or that the carrier is minimal. The oracle condition is especially diagnostic: a large remaining failure rate means branch search or proof search, not carrier invention, is still the main bottleneck.

\Q{R9.2}{What makes model-assisted search deterministic and auditable?}
\ans Freeze the proposal stream for search replay, restrict the proposal language, and validate the final term independently. Temperature zero is not, by itself, a cross-platform reproducibility guarantee. Search replay should not require regenerating identical model outputs.

Accept model output only as a small typed proposal language: an allowed carrier expression, a helper signature, a component ordering, or a recursion-scheme choice. Do not execute arbitrary suggested tactics, imports, attributes, commands, or foreign code. Such code can affect the elaboration environment and the meaning of later statements, not just propose a harmless term. Lean's proof-validation guidance makes the distinction between a valid proof and the intended statement explicit \cite{leanvalidation}.

A receipt should record the complete proposal bytes, their canonical parsed form, all rejection reasons, the model and tokenizer identifiers, prompt/template version, decoding settings, and a hash of the allowed environment and grammar. To reproduce the mathematical result, store the closed program and proof. To reproduce the search, replay the accepted/rejected proposal stream and logical scheduling events. Reproducing the model's own numerical inference additionally requires its weights, backend, and numerical configuration and may be a stronger, separate service.

\textbf{Proposal scope.} An ordering suggestion within $G_B$ changes only discovery order. A carrier expression not already represented in $G_B$ changes the search space. The latter may be valuable, but label it an extension $G_B\subseteq G'_B$ and record which results used it. A model failure or absence of a proposal cannot justify a negative result about either grammar. Maintain the symbolic baseline lane.

\textbf{Four audit levels.} A proof receipt supports independent checking of the frozen statement and axiom policy. A realization receipt supports checking the executable program's equivalence to the logical term. A search receipt supports replay of the explored bounded grammar and scheduling decisions. A model receipt supports attribution of proposals to the recorded inference configuration. Keeping these levels separate prevents ``kernel checked'' from being misread as ``search exhaustive,'' ``compiled behavior independently proved,'' or ``model deterministic.''

The model remains outside the mathematical trusted base when proposals are fully validated and cannot alter the checker, environment, or original statement. It can still damage usefulness by wasting time, biasing ranking, or hiding alternatives. Those are addressed by accounting, replay, grammar tracking, and fair fallback, not by proof checking alone.
